% !TEX root = ../tex/3-6.tex
% 第3章：本文は青、補足は緑
\newcommand{\asmain}{%
    \setbeamercolor{frametitle}{fg=white,bg=main}%
    \setbeamercolor{structure}{fg=main}%
    \setbeamercolor{itemize item}{fg=main}%
    \setbeamercolor{itemize subitem}{fg=main}%
    \setbeamercolor{enumerate item}{fg=main}%
}
\newcommand{\assuppl}{%
    \setbeamercolor{frametitle}{fg=white,bg=examplecol}%
    \setbeamercolor{structure}{fg=examplecol}%
    \setbeamercolor{itemize item}{fg=examplecol}%
    \setbeamercolor{itemize subitem}{fg=examplecol}%
    \setbeamercolor{enumerate item}{fg=examplecol}%
}
\asmain


%=============================================
\begin{frame}{3.6 Closing Remarks}
    \begin{block}{この章でたどった第二の道}
        変分（第2章）とは別に、
        EBMに根ざすスコア視点から拡散モデルへ至る道を見た。
    \end{block}
    \begin{itemize}
        \item 出発点：分配関数が扱いにくい
        \item 道具：スコア$\nabla_{\bm{x}}\log p(\bm{x})$
        \item 到達点：複数ノイズのNCSNと、DDPMとの対応
    \end{itemize}
\end{frame}
%=============================================
\asmain
\begin{frame}{第3章の道筋（前半）}
    \begin{enumerate}
        \item \textbf{EBM}
        エネルギーで分布を書く。最尤は$Z_\phi$で行き詰まる。
        \item \textbf{スコア}
        $s=\nabla_{\bm{x}}\log p$なら正規化定数が消える。
        \item \textbf{スコアマッチング}
        真のスコアなしで場を合わせるが、トレース計算が残る。
    \end{enumerate}
\end{frame}
%=============================================
\asmain
\begin{frame}{第3章の道筋（後半）}
    \begin{enumerate}
        \setcounter{enumi}{3}
        \item \textbf{DSM}
        ノイズを入れ、条件付きスコアを教師にする。\\
        Tweedieにより、スコア学習＝ノイズ除去学習になる。
        \item \textbf{NCSN}
        複数ノイズで条件付きスコアを学び、
        アニーリング・ランジュバンで生成する。
        \item \textbf{DDPMとの比較}
        公理と手続きは違うが、同じ梯子の別の書き方である。
    \end{enumerate}
\end{frame}
%=============================================
\asmain
\begin{frame}{一本の流れ}
    \begin{block}{EBM $\to$ スコア $\to$ DSM $\to$ NCSN}
        分配関数を避ける
        $\;\to\;$
        スコアを直接扱う
        $\;\to\;$
        ノイズで教師を作る
        $\;\to\;$
        複数スケールで粗から細へ戻す
    \end{block}
    \begin{block}{核心}
        密度そのものではなく対数密度の勾配を学び、
        確率的な流れでデータへ戻す。
    \end{block}
\end{frame}
%=============================================
\asmain
\begin{frame}{離散の梯子から連続へ}
    NCSNのノイズ列もDDPMのステップも、有限個の段である。\\
    次章では段のあいだを埋め、ノイズレベルを時間の連続関数にする。
    \[
        p_{\mathrm{data}}
        \xrightarrow{\;\text{forward SDE}\;}
        p_T\approx p_{\mathrm{prior}},
        \qquad
        p_{\mathrm{prior}}
        \xrightarrow{\;\text{reverse dynamics}\;}
        p_{\mathrm{data}}
    \]
    \begin{itemize}
        \item 前向きSDE：データを徐々にノイズへ崩す
        \item 逆ダイナミクス：時間依存スコア$\nabla_{\bm{x}}\log p_t(\bm{x})$で戻す
    \end{itemize}
\end{frame}
%=============================================
\asmain
\begin{frame}{第4章で行うこと}
    \begin{enumerate}
        \item \textbf{連続時間へ}
        DDPMとNCSNを、一つのSDEの異なる離散化として統一する
        （Score SDE）。
        \item \textbf{二つの視点をつなぐ}
        変分視点とスコア視点を、微分方程式を解く生成としてつなぐ。
    \end{enumerate}
    \begin{block}{ねらい}
        理論の見通しを良くし、
        遅いサンプリングに対する高度な数値法への道を開く。
    \end{block}
\end{frame}
%=============================================
\assuppl
\begin{frame}{\supplmark：この章の到達目標との対応}
    章頭の「できるようになること」への対応である。
    \begin{itemize}
        \item スコアとランジュバンによる生成 $\;\to\;$ 3.1–3.2
        \item 単一ノイズのDSMとTweedie $\;\to\;$ 3.3
        \item 複数ノイズのNCSNとDDPMノイズ予測との関係 $\;\to\;$ 3.4–3.5
    \end{itemize}
\end{frame}
%=============================================
\asmain
\begin{frame}{第3章のまとめ}
    \begin{itemize}
        \item スコア視点は、EBMの分配関数問題を迂回する第二の入口である
        \item DSMとTweedieが、学習を単純な回帰とノイズ除去へ落とす
        \item NCSNは複数スケールへ拡張し、DDPMと同じ情報を別パラメータで持つ
        \item 残る課題は密な離散化による遅さであり、第4章の連続時間へ進む
    \end{itemize}
\end{frame}
